Estimating the covariance matrix of asset returns is a fundamental first step in portfolio optimization. Many of the results that follow, from the mean-variance efficient frontier to the optimal portfolio weights, hinge on how well this matrix is estimated. Forecasting the covariance matrix of asset returns conditional on past information is another essential task in portfolio optimization for hedging and re-balancing portfolio holdings. Much attention has been devoted to conditional forecasting of the covariance matrix of financial returns, especially after the introduction of the generalized autoregressive conditionally heteroskedastic (GARCH) models by \textcite{ENGLE1982} and \textcite{BOLLERSLEV1986}. These models have subsequently been expanded to the multivariate case, allowing for richer dynamics by accounting also for the covariances between the return series, giving rise to multivariate GARCH models. Estimation of conventional multivariate GARCH models, however, is often not feasible due to a lack of degrees of freedom. The covariance matrix of a portfolio composed of $N$ distinct assets has $N(N+1)/2$ parameters, while the number of available observations from a sample of size $T$ is only $N \times T$. Therefore, when the number of assets is of the same order of magnitude as the sample size, that is, when $N$ and $T$ are of the same order of magnitude, the covariance matrix has $O(N^2)$ parameters that need to be estimated from only $O(N^2)$ noisy observations, which is statistically infeasible. A variant of the multivariate GARCH models, the conditional correlation class of models including the constant conditional correlation (CCC) model introduced by \textcite{BOLLERSLEV1990} and the dynamic conditional correlation (DCC) model introduced by \textcite{ENGLE2002}, alleviates this problem by dividing the estimation process into two steps. In the first step, univariate GARCH models are estimated for each of the return series and then the devolatilized return series are obtained. Next, in the second step, taking the univariate GARCH parameter estimates as given, the conditional correlation matrix of the return series is estimated. Given that the number of parameters to be estimated in the conditional correlation process is independent of the number of series to be correlated, these models can allow for the estimation of potentially large correlation matrices.

Estimation of the conditional correlation models can still be challenging in large dimensions. There are two main sources of difficulty. The first is that the models require the estimation of the unconditional correlation matrix, for which the usual estimator, the sample correlation matrix, is not feasible in large dimensions. The second problem is that the maximum likelihood estimation of the dynamic correlation parameters requires inverting the conditional covariance matrix, which could make the likelihood function numerically unstable due to the conditional covariance matrix being possibly ill-conditioned. We now turn to each issue separately.

The sample covariance matrix, or the sample correlation matrix for standardized variables, is a very poor estimator in large dimensions. While under standard asymptotics the sample covariance matrix is well-conditioned and asymptotically efficient as an estimator for the population covariance matrix, when the dimension of the matrix is large compared to the sample size, this matrix is often ill-conditioned and inconsistent. What is more, the standard assumptions underlying standard asymptotics are no longer appropriate when the dimension of the matrix is not negligible with respect to the sample size. Therefore, we need to adopt a new asymptotic framework in which the number of parameters is allowed to go to infinity together with the sample size. In this situation, the sample covariance matrix tends to be far from the population covariance matrix even asymptotically, and its spectrum converges almost surely to a limit law that is related to, but fundamentally different from, the spectrum of the population covariance matrix. The appropriate asymptotic framework to employ, therefore, is large dimensional asymptotics where both the dimension of the covariance matrix $N$ and the sample size $T$ go to infinity, with their ratio converging to a non-zero limit called the concentration ratio $C=N/T > 0$. The asymptotic spectral distribution of the sample covariance matrix is governed by the Marčenko-Pastur equation and is related to the spectral distribution of the population covariance matrix through the Stieltjes transform. This fundamental result was first established by \textcite{MP1967} for a general class of random matrix ensembles, and later elaborated by \textcite{SILVERSTEIN1995} and \textcite{SILVERSTEIN_CHOI_1995} for the special case of random matrix ensembles of the sample covariance type. Basic analytical results in the random matrix theory (RMT) literature indicate that the eigenvalues of the sample covariance matrix are biased, with the largest eigenvalues significantly biased upwards and the smallest ones biased downwards. Therefore, it would be reasonable to correct for the in-sample bias of the sample eigenvalues by replacing them with a shrinkage estimator that pulls down the largest eigenvalues and pushes up the smallest ones. To this end, a first approximation would be to shrink the sample covariance matrix toward a structured shrinkage target matrix. This is the approach of \textcite{LEDOITWOLF2004} in which they shrink the sample covariance matrix toward a multiple of the identity matrix by computing the optimal convex linear combination of the sample covariance matrix and the identity matrix. In their estimation scheme, \textcite{LEDOITWOLF2004} estimate the optimal linear shrinkage intensity under the Frobenius norm of the difference with the population covariance matrix, and show that their linear shrinkage estimator is always well-conditioned and more efficient than the sample covariance matrix. In their method, shrinkage takes place at the level of the eigenvalues, where they are all drawn toward a central target value, hence smoothing the spectrum of the sample covariance matrix and reducing estimation error by minimizing the noise involved in the sample eigenvalues. Although the linear shrinkage estimator of \textcite{LEDOITWOLF2004} is the optimal convex linear combination of the sample covariance matrix and the identity matrix, this estimator is not optimal in the more general sense when considering the quadratic loss under the Frobenius norm of the difference with the population covariance matrix, due to the non-linear nature of the problem. In linear shrinkage, all sample eigenvalues are shrunk with the same intensity, regardless of their magnitude, whereas in general, under the asymptotic relationship between the sample and population eigenvalues, a different shrinkage intensity has to be applied to each sample eigenvalue, and hence, linear shrinkage is only a first approximation to the problem. \textcite{LEDOITPECHE2011} develop an \textit{oracle} shrinkage formula for non-linear shrinkage intensities that is optimal under the Frobenius norm for the more general class of rotation-equivariant estimators. This \textit{oracle} shrinkage formula is the basis for the non-linear shrinkage estimator of \textcite{LEDOITWOLF2015}. Their estimation methodology is based on the inversion of the Marčenko-Pastur equation, from which they obtain an estimate for each individual population eigenvalue. They achieve this by the numerical inversion of what they call the QuEST function, where QuEST stands for \textit{Qauntized Eigenvalues Sampling Transform} and that it is a multivariate non-random function that discretizes the Marčenko-Pastur equation. They further show that their estimates of the population eigenvalues are optimal and consistent under large dimensional asymptotics in the sense that the mean squared deviation between the estimator and the population eigenvalues converges almost surely to zero. Substituting the eigenvalues of the sample covariance matrix with the estimates of the population eigenvalues obtained via non-linear shrinkage yields a covariance matrix estimator that is asymptotically optimal for the population covariance matrix among the class of rotation-equivariant estimators. That is to say, the Frobenius norm of the difference between the estimator and the population covariance matrix, normalized by the dimension of the covariance matrix, is minimized over the set of all covariance matrix estimators that are rotation-equivariant. In addition, just like the linear shrinkage estimator, the non-linear shrinkage estimator is always well-conditioned.

Besides the shrinkage estimation methodology discussed above, there exist other methods for obtaining improved estimators for the covariance matrix. One such method involves incorporating additional prior information in the estimation process, such as sparseness or factor structure. These methods often require imposing stricter assumptions on the models and can even lead to consistent covariance matrix estimators if the assumptions are satisfied. However, such additional prior information might not always be available or trustworthy. Therefore, in the general case and in the absence of reliable prior information, it is reasonable to resort to a general class of estimators that are equivariant under orthogonal rotations of the sample data. This means that if the data were rotated by any orthogonal matrix, the estimator would also rotate by the same matrix and in exactly the same way. Estimators of this kind are generally called rotation-equivariant estimators and share the same set of eigenvectors. For the class to which the sample covariance matrix belongs, this means that all estimators have the same set of eigenvectors as the sample covariance matrix, and, as has been established earlier, the non-linear shrinkage estimator of \textcite{LEDOITWOLF2015} is the asymptotically optimal estimator within the class. The same eigenvalue shrinkage approach can be applied to the sample correlation matrix as well. This is particularly useful for estimating the unconditional correlation matrix in the conditional correlation models, and it is the approach that we take in this thesis for correlation targeting in the conditional correlation models. Hence, by means of non-linear shrinkage of the sample correlation matrix, the first of the two fundamental problems regarding estimation of conditional covariance matrix models in large dimensions is addressed.

The second salient problem regarding the estimation of covariance matrix models in large dimensions is related to the maximum likelihood estimator being numerically unstable due to the conditional covariance matrix being possibly ill-conditioned. Computing the likelihood function requires inverting an $N \times N$ conditional covariance matrix $T$ times throughout the estimation sample. Other than this operation being computationally burdensome, as it requires $O(N^3)$ operations for each inversion of the $N \times N$ conditional covariance matrix, the maximization problem is ill-posed, as the likelihood function involves the inverse of large dimensional conditional covariance matrices that are ill-conditioned and can hence dramatically increase the estimation error in the maximum likelihood estimates. Furthermore, the maximum likelihood estimator of the unconditional covariance (correlation) matrix coincides with the sample covariance (correlation) matrix under the assumption of normality, which, as has been discussed above is substantially biased. Using a well-conditioned estimator such as the shrinkage estimators discussed above in place of the sample covariance matrix can partly alleviate the problem related to the biasedness of the sample covariance matrix in the model. The issues regarding the computation cost and the ill-posedness of the likelihood function, however, remain outstanding. One particular way to handle this problem is to consider the joint likelihoods of pairs of assets instead of the full likelihood function and then take the average of their logarithms. This approach is known as the composite likelihood method and is computationally feasible, and it avoids the ill-posedness of the full information maximum likelihood estimation since the conditional covariance matrix in the joint likelihood function of each pair of assets is only $2$-dimensional, regardless of the total number of assets being considered in the model. The drawback of the composite likelihood method, however, is that it incurs an efficiency cost, as the potential dependence between the bivariate likelihoods is ignored. The composite likelihood method was extensively studied by \textcite{PAKELETAL} who derived its asymptotic properties. They show that under certain regularity conditions, the estimates obtained from maximizing the composite likelihood function are asymptotically normal, and that in spite of being inefficient, they are nevertheless consistent. The construction of the composite likelihood function as outlined in \textcite{PAKELETAL} can proceed in two ways. The first, dubbed the 2MCLE (2-Maximum Composite Likelihood Estimator) -- where the  "2" represents the two-step structure where the first step estimates the unconditional covariance (correlation) matrix, e.g., via shrinkage of the sample covariance (correlation) matrix, and the second step estimates the parameters of the conditional correlation model -- is based on all bivariate combinations of the $N$ assets which contains $\binom{N}{2}$ pairs. The second method, called the 2MSCLE, is based on the contiguous pairs of assets which includes only $N-1$ pairs of assets. In their remarks, they note that the difference between the 2MCLE and 2MSCLE is minuscule for large $N$, and hence in the estimation of the models we will make use of the 2MSCLE for it is more scalable than the 2MCLE.

The estimation mechanisms discussed above are aimed at obtaining improved covariance matrix estimators that are well-conditioned and robust in large dimensions. Indeed, the mean-variance optimization principle which has laid the foundations of modern portfolio theory and is the core principle in portfolio optimization, hinges crucially on the statistical estimation scheme that precedes it. In the original formulation of the portfolio choice problem by \textcite{MARKOWITZ}, the optimization of the portfolio holdings is divided into two stages. Of these two stages, the seminal work of \textcite{MARKOWITZ} is only concerned with the second stage which involves finding the optimal portfolio weights in an investment universe where the expected future returns and covariances of the securities are known. Needless to say, the optimal portfolio weights that are subsequently obtained from the optimization process are only optimal insofar as the estimated expected returns and covariances that led to them are reliable estimates of the true population parameters. When the number of securities for constructing the portfolio are small to moderate, traditional estimation schemes often provide reliable estimates that can be used in the optimization process. However, as the number of the securities in the investment universe becomes large, say of the same order of magnitude as the sample size, the number of free parameters in the models becomes too large to be consistently estimated from the available data. This is a symptom commonly known as the curse of dimensionality, and it can severely limit the breadth and scope of estimation analysis in large dimensional data sets. Unless the limits posed by the large dimensional nature of the data are addressed, it is not possible to obtain reliable estimates of the parameters that are required for the optimization process. This can have a significant toll on the portfolio choice problem in general, where if the estimates obtained in the first stage are not reliable, any decision based on them is bound to be wrong and misinformed. The real challenge in practical applications is that the optimization process can significantly amplify the error involved in the estimated parameters, and if the estimation error is large to begin with, the final outcome would be too erroneous to have any value to be employed in a practical setting. This was famously pointed out by \textcite{MICHAUD}, who characterized portfolio optimization as "error maximization." The estimation methods that have been introduced above aim at alleviating exactly this problem. Both methods have been extensively studied in the literature, and the results are promising. Combining the two methods can provide yet stronger results that can handle potentially very large portfolios composed of thousands of assets. \textcite{ELW2019} were the first to combine these two methods. They show that the best way to apply the shrinkage method is by shrinking the sample correlation matrix for estimation of the unconditional correlation matrix rather than the direct application of shrinkage to the estimated conditional covariance matrix. Next, they use the composite likelihood method to estimate the dynamic parameters of the DCC model. Their results are promising. The combined shrinkage and composite likelihood methods performs desirably well both in simulations and in out-of-sample evaluations in applications to real data, even and especially when the number of assets under consideration is very large, in the order of a thousand. Our approach in this thesis is founded on the same basis. We study the properties of the DCC model estimated via shrinkage (both linear and non-linear) and composite likelihood over Monte Carlo simulations and in applications to real data. \textbf{This section for results once we have them later.}

Finally, the only remaining concern would be the vector of expected returns. Mean-variance optimization requires estimates of both the covariance matrix and the vector of expected returns of the assets. Estimation of these two parameters, however, need not necessarily follow the exact same statistical framework. As long as estimators for both quantities are reliable and consistent, they can be used for solving the portfolio optimization problem. Numerous models have been proposed for estimating the vector of expected returns. The most prominent examples are the capital asset pricing model of \textcite{SHARPE1964} and \textcite{LINTNER1965}, the factor pricing model of \textcite{ROSS1976341} and the factor models of \textcite{FAMAFRENCH1992} and \textcite{FAMAFRENCH2015}. See \textcite{GREENETAL} for a comprehensive review and the references therein. The literature on estimating expected returns is vast and beyond the scope of this thesis. Regardless, the covariance matrix models studied here may be used without any limits together with any model for the expected returns for optimization.